% Root construction C.3ag; independent full audits by frobenius and real_weil.
\paragraph{C.3ag (Refuted): ordinary restriction in dimension six.}
We specify an excellent normal strictly henselian local
complex ring $R$ of dimension six, with regular puncture
$X$ and regular punctured principal divisor $Y$, and a
surjection
\[
 \operatorname{coker}\bigl(
 \mathrm{CH}^2(X)_{\mathrm{alg}}\longrightarrow
 \mathrm{CH}^2(Y)_{\mathrm{alg}}\bigr)
       \longrightarrow \operatorname{Pic}^0(C)(\mathbb C),
 \qquad C=V(z_0^4+z_1^4+z_2^4)\subset\mathbb P^2.
\]
The algebraic-equivalence groups use actual product
families. All differential sheaves below are relative
to $\mathbb C$. The surface in this construction is a
specified very general member with a Hodge-generic
complex choice; unlike C.3ae, we do not assert that it
is the Fermat quartic or that its coefficients are integers.

\emph{1. The surface and the precise genericity in the data.}
Fix the line $D=V(w_0,w_1)\subset\mathbb P^3$. Let
$\mathcal U$ be the open set of smooth quartics in
\[
 \left\{\sum_{\substack{|e|=4\\e_0+e_1>0}}
                 u_e w^e=0\right\}.
\]
There are thirty coefficients. We choose a Hodge-generic
complex point in this family for which
\[
 \operatorname{NS}(W)=\mathbb Z H_W\oplus\mathbb Z[D],
 \qquad \operatorname{End}_{\mathrm{Hdg}}(T_W)=\mathbb Q,
 \qquad T_W=\operatorname{NS}(W)_{\mathbb Q}^{\perp}
                  \subset H^2(W,\mathbb Q).
\]
Here ``Hodge-generic'' is part of the choice, not a
claim about every embedding of its coefficient field.
The existence, exact source inputs, and the consequence
for mixed cycle classes are supplied in the following
lemma. In particular they are not additional conjectures.

% Auxiliary input for the parent's prospective dimension-six construction.
% No complete local-geometric assertion is made in this fragment.
\paragraph{Auxiliary Hodge calculation for a quartic containing a line.}
Fix a line $D_0\subset\mathbb P^3_{\mathbb C}$ and let
$B\subset\mathbb P H^0(\mathcal I_{D_0}(4))$ be the nonempty open subset
parametrizing smooth quartics. We choose a \emph{Hodge-generic complex
point} $b\in B$ as part of the data, in the precise sense described below.
Write $W=W_b$, $H=c_1(\mathcal O_W(1))$, and $D=D_0\subset W$.
We can and do choose $b$ so that
\[
 \begin{gathered}
 \operatorname{NS}(W)=\mathbb ZH\oplus\mathbb ZD,
 \qquad
 \begin{pmatrix}H^2&H\cdot D\\H\cdot D&D^2\end{pmatrix}
 =\begin{pmatrix}4&1\\1&-2\end{pmatrix},\\
 T=\operatorname{NS}(W)_{\mathbb Q}^{\perp}
          \subset H^2(W,\mathbb Q),
 \qquad \operatorname{End}_{\mathrm{Hdg}}(T)=\mathbb Q.
 \end{gathered}
\]
Let $C$ be the Fermat plane quartic and put $S=W\times C$.
Then the largest rational sub-Hodge structure of $H^3(S,\mathbb Q)$
whose Hodge types are contained in $\{(2,1),(1,2)\}$ is
\begin{equation}\label{eq:N6-level-one}
       \operatorname{NS}(W)_{\mathbb Q}\otimes H^1(C,\mathbb Q).
\end{equation}
In particular, for every smooth projective complex curve $P$ and every
$z\in\mathrm{CH}^2(W\times C\times P)$, the component of its rational
cycle class in $H^3(S,\mathbb Q)\otimes H^1(P,\mathbb Q)$ belongs to
\[
 \operatorname{NS}(W)_{\mathbb Q}\otimes
 H^1(C,\mathbb Q)\otimes H^1(P,\mathbb Q).
\]

\emph{Choice of the quartic and its Hodge-generic complex point.}\par
Theorem~1.1 and Corollary~1.3(c) of
\citet[p.~2 of the opened author version]{frBrevikNollet2011}
apply to $D_0$ with $d=4$: $\mathcal I_{D_0}(3)$ is globally generated
and $H^0(\mathcal I_{D_0}(2))\ne0$. They give the asserted Picard group
for a very general member of $B$. A smooth quartic is a K3 surface,
and adjunction gives the displayed intersection matrix.

For completeness, we also justify the endomorphism condition rather
than inferring it solely from the Picard rank. Locally mark the second
cohomology. The Local Torelli Theorem and its line-bundle formulation
in \citet[Proposition~6.2.8 and Section~6.2.4, pp.~107--108 of the
opened author version]{frHuybrechtsK3Book} identify the local deformations
preserving $H,D$ with an open subset of the period quadric in
$(\mathbb QH+\mathbb QD)^\perp_{\mathbb C}$. It has dimension $18$.
These are locally realized by quartics containing a line: very ampleness
of $H$ persists; Riemann--Roch, $D^2=-2$, and $H\cdot D=1$ make $D$
effective, and an effective divisor of hyperplane degree one is a line.
A projective coordinate change puts that line back at $D_0$.

For every nonscalar $a\in\operatorname{End}_{\mathbb Q}(T)$, the
condition that $a$ is a Hodge endomorphism implies that the period line
is an eigenline of $a$. The eigenlines form a finite union of proper
linear subspaces. They cannot contain the nondegenerate irreducible
period quadric, which spans the ambient projective space. Thus each
such Hodge-endomorphism locus is proper. There are countably many
rational endomorphisms. Applied to the polarized weight-zero variation
$\operatorname{End}(T)$, the Hodge-locus theorem of
\citet[Theorem~1.1 and Corollary~1.3, p.~2]{frCDK1994}
shows that these loci are algebraic (with their monodromy translates).
Avoiding their countable union, together with the extra Picard loci,
produces the stipulated complex point. This is the meaning of the
Hodge-generic choice. We do not assert that every embedding of a field
with coefficientwise algebraically independent quartic parameters has
this property.

\emph{The product Hodge group and its representations.}
By \citet[Theorem~2.2.1 and Example~2.2.3, p.~211]{frZarhin1983},
\[
                    \operatorname{Hg}(T)=\operatorname{SO}(T).
\]
The exact original theorem and its proof in Section~2.5, pp.~212--213,
have been inspected. By \citet[Lemma~2, p.~106, and Example~3,
pp.~106--107]{frBeauville2014}, the Jacobian of $C$ is isogenous to
$E^3$, where $E=\mathbb C/\mathbb Z[i]$. Hence, for $V=H^1(E,\mathbb Q)$,
we have $H^1(C,\mathbb Q)\simeq V^{\oplus3}$ and
$\operatorname{Hg}(V)=\operatorname{Res}^{1}_{\mathbb Q(i)/\mathbb Q}
\mathbb G_m=:U$.

Let $G=\operatorname{Hg}(T\oplus V)$. Its two projections are surjective,
and $G$ is connected. The image of $G^{\mathrm{der}}$ on
$\operatorname{SO}(T)$ is the full group, while its image on the torus
$U$ is trivial. Therefore $\operatorname{SO}(T)\times1\subset G$;
surjectivity on $U$ then gives
\[
                   G=\operatorname{SO}(T)\times U.
\]
This argument uses only that the special orthogonal group in rank $20$
is connected and semisimple; no independence assumption on periods
of $W$ and $E$ is being inserted.

Over $\mathbb Q(i)$, the representation $T\otimes V$ is the direct sum
of the two twists of the standard $\operatorname{SO}(T)$ representation
by the two opposite nontrivial characters of $U$. Both summands are
absolutely irreducible, they are nonisomorphic, and complex conjugation
exchanges them. A rational invariant subspace therefore contains both
or neither. Thus $T\otimes V$ is a simple rational Hodge structure.
Moreover
\[
                 (T\otimes V)^{3,0}=T^{2,0}\otimes V^{1,0}\ne0.
\]
It consequently has level three. Semisimplicity of polarizable rational
Hodge structures implies that every nonzero sub-Hodge structure of
$(T\otimes V)^{\oplus3}$ has a simple constituent isomorphic to
$T\otimes V$; in particular none has level at most one.

Since $H^1(W)=H^3(W)=0$, Kunneth gives
\[
 H^3(S,\mathbb Q)=
 \bigl(\operatorname{NS}(W)_{\mathbb Q}\otimes H^1(C,\mathbb Q)\bigr)
 \oplus\bigl(T\otimes H^1(C,\mathbb Q)\bigr).
\]
The first summand has exactly the two indicated Hodge types and the
second has no nonzero sub-Hodge structure with those types. This proves
\eqref{eq:N6-level-one}.

Finally, the indicated Kunneth component of the class of $z$ is a
rational Hodge class of type $(2,2)$. Poincare duality on $P$ turns it
into a morphism of Hodge structures
$H^1(P,\mathbb Q)(-1)\to H^3(S,\mathbb Q)$. The source has types
$(2,1),(1,2)$, so its image lies in \eqref{eq:N6-level-one}, proving
the claimed consequence for every such cycle and parameter curve.


Let $K_0\subset\mathbb C$ be the field generated over
$\overline{\mathbb Q}$ by these thirty coefficients.
Choose thirty-six elements $a_{ijk}$, indexed by
$0\le i\le2$, $0\le j\le3$, $0\le k\le2$, algebraically
independent over $K_0$. Put
\[
 S=W\times C,\quad M=\mathcal O_W(1)\boxtimes K_C,
 \quad a=c_1(M),\quad
 s_i=\sum_{j,k}a_{ijk}w_jz_k.
\]
This specifies the geometric generic ordered triple
of sections over the chosen surface, with any complex
embedding extending the stated data. We prove the
two required open conditions for this triple below.

\emph{2. The section geometry and the dimension-six germ.}
The very ample bundle $M$ embeds the smooth threefold
$S$. Three general hyperplanes in this embedding have
transverse zero-dimensional intersection. This follows
successively from Bertini: after each smooth section,
the remaining very ample system is still without base
points. The last intersection is reduced, and its length is
\[
                        \int_S a^3
              =3\left(\int_W H_W^2\right)\deg K_C=48.
\]
The resulting nonempty Zariski-open condition is
defined over $K_0$, so it holds for the chosen triple.
Write $B=V(s_0,s_1,s_2)\subset S$ for these forty-eight
points, and put
\[
 V=\mathbb P^2_x\times S,\qquad
 L=\mathcal O_{\mathbb P^2}(1)\boxtimes M,
 \qquad \ell=\sum_{i=0}^2 x_i s_i,
 \qquad H=V(\ell)\subset V.
\]
The projection $p:H\to S$ is a $\mathbb P^1$-bundle
over $U=S\setminus B$, and has fibre $\mathbb P^2$
over each point of $B$. The latter locus has codimension
two in $H$. At a point of $B$, the three differentials
$ds_i$ form a basis of the cotangent space of $S$.
For every nonzero $x$, the differential
$\sum x_i ds_i$ is therefore nonzero. Together with
the bundle description over $U$, this proves that
$H$ is smooth. It is integral: its integral open
$\mathbb P^1$-bundle is dense and the exceptional
fibres have too small a dimension to be components.

Set
\[
 B'=\mathbb P^2_x\times W\times\mathbb P^2_z,
 \quad\mathcal N=\mathcal O_{B'}(-1,-1,-1),\quad
 G=z_0^2z_1^2+z_0^2z_2^2+z_1^2z_2^2.
\]
Inside $\operatorname{Tot}(\mathcal N)$ define
$Z=V(F+t\ell G)$, where $F=z_0^4+z_1^4+z_2^4$.
Contract by the Segre coordinates $y_{ijk}=tx_iw_jz_k$,
take the reduced closed image, its finite normalization,
and the strict henselization at the unique point over
the vertex. This is $R$. All these operations and all
their equations have now been specified. The proof of
C.3ae applies without a base-point-free hypothesis on
the $s_i$: $Z$ is integral, smooth near its connected
exceptional divisor $V$, and its regular projective
model $\mathcal Z$ over $R$ has
$\mathfrak m_R\mathcal O_{\mathcal Z}=\mathcal I_V$.
Its complement is the regular puncture $X$.
The dimension is $\dim V+1=6$.

The function $f=\sum a_{ijk}y_{ijk}$ pulls back to
$t\ell$. Its strict transform is the total space of
$L_H^{-1}$ on $H$, before the indicated base changes.
Write $\mathcal D$ for the corresponding regular
divisor on $\mathcal Z$. Then
$Y=\mathcal D\setminus H$ is the regular punctured
principal divisor required in the statement, and
$\mathcal D$ meets $V$ transversely in $H$.

\emph{3. An integral incidence invariant and its boundary.}
Put $h=c_1(\mathcal O_{\mathbb P^2}(1))|_H$ and
$\tau=4[D]-H_W$. Adjunction gives
\[
 H_W^2=4,\quad H_W[D]=1,\quad [D]^2=-2,
 \quad \deg(\tau H_W)=0,\quad\deg(\tau[D])=-9.
\]
For a smooth projective connected parameter $P$,
including a point, localization and self-intersection
give
\[
 \sigma_P:\mathrm{CH}^2(Y\times P)\longrightarrow
 \frac{\mathrm{CH}^2(H\times P)}
             {(h+a)\mathrm{CH}^1(H\times P)}.
\]
Indeed we extend a cycle to $\mathcal D\times P$ and
restrict to $H\times P$; a different extension changes
the restriction by $-(h+a)$ times a divisor class.

We need only a divisor decomposition, not a
projective-bundle description of all of $H$. Removing
$p^{-1}(B)\times P$ changes no Picard group, because
it has codimension two in a regular scheme. Removing
$B\times P$ similarly changes no Picard group of
$S\times P$. On the complement $H_U\times P$ is a
$\mathbb P^1$-bundle, with the line bundle of class
$h$ of relative degree one. Therefore, integrally,
\[
 \mathrm{CH}^1(H\times P)=
 p^*\mathrm{CH}^1(S\times P)
       \oplus h\,p^*\mathrm{CH}^0(S\times P).
\]
These equalities can also be read directly from Weil
divisors: deletion in codimension at least two preserves
divisor classes on a regular scheme, and the generic
fibre degree splits the divisor group of a projective
bundle. The only pullbacks by the possibly nonflat
map $p$ used here are of divisor classes or of their
products.

The hypersurface class $[H]=h+a$ in
$\mathbb P^2\times S$ gives the exact pushforwards
\[
            p_*1=0,\quad p_*h=1,\quad p_*h^2=a,
                      \quad p_*h^3=0.
\]
For $\alpha\in\mathrm{CH}^2(H\times P)$ set
\[
 \Phi_P(\alpha)=\operatorname{pr}_{C\times P,*}
          \left(\tau\,p_*\bigl((h-2a)\alpha\bigr)\right)
          \in\operatorname{Pic}(C\times P).
\]
Write a divisor on $H\times P$ as $\gamma_0+h\gamma_1$
by the preceding decomposition. Projection formula gives
\[
 p_*\bigl((h-2a)(h+a)(\gamma_0+h\gamma_1)\bigr)
                         =-3a^2\gamma_1.
\]
Its product with $\tau$ has zero pushforward to
$C\times P$: the $H_W^2$ term has codimension three
on $W$, and the $2H_WK_C$ term has coefficient
$2\deg(\tau H_W)=0$. Thus $\Phi_P$ descends to the
quotient defining $\sigma_P$, and
$\rho_P=\Phi_P\sigma_P$ is well defined integrally.
Localization, refined restriction, products and proper
pushforward commute with taking a parameter fibre.
Hence these $\rho_P$ are compatible with actual families.

If a cycle on $X\times P$ is closed on
$\mathcal Z\times P$ and restricted to $V\times P$,
write its class as $b=b_0+h b_1+h^2b_2$ using the
projective bundle $V=\mathbb P^2\times S$. Transverse
restriction to $H$, followed by the preceding identities,
gives
\[
 \rho_P(\Gamma|_{Y\times P})=
 \operatorname{pr}_{C\times P,*}
                  \bigl(\tau(b_0-a b_1-2a^2b_2)\bigr).
\]
This formula is independent of the closure. It has
been derived on the fixed scheme times $P$, without
any assertion about spreading a varying etale model.

\emph{4. Why three sections suffice on this surface.}
Let $F_W$ be the equation of $W$, and set
\[
 \mathcal R=\mathbb C[w_0,w_1,w_2,w_3]/
                 (\partial F_W/\partial w_0,\ldots,
                  \partial F_W/\partial w_3).
\]
The Euler and normal sequences used in C.3ae work
for every smooth quartic. They give
\[
 \frac{H^1(W,\Omega_W^1)}{\mathbb C H_W}
        \simeq\mathcal R_4^*,\qquad
 H^1(W,\Omega_W^1(1))\simeq\mathcal R_3^*,
 \qquad \dim\mathcal R_4=19,\quad\dim\mathcal R_3=16.
\]
Under these identifications, multiplication by $w_j$
is dual to $w_j:\mathcal R_3\to\mathcal R_4$ and kills
$H_W$. The sum
$\mathcal R_3^{\oplus4}\to\mathcal R_4$ is surjective:
every polynomial of degree four is a sum of $w_j$
times polynomials of degree three, and this property
survives the homogeneous quotient. Consequently the
four dual multiplication maps have zero common kernel
modulo $\mathbb C H_W$.

The class $\tau=4[D]-H_W$ is nonzero modulo $H_W$;
indeed $\tau H_W=0$ in degree and $\tau^2=-36$.
Thus some linear form $w$ satisfies
$\mu_w(\tau)\ne0$ in $H^1(W,\Omega_W^1(1))$.
The necessary rank condition on our triple is the
injectivity of
\[
 \mathbb C\tau\otimes H^0(C,K_C)
   \longrightarrow
 H^0(\mathbb P^2,\mathcal O(1))
                 \otimes H^1(W,\Omega_W^1(1)),
 \quad
 \tau\otimes v\longmapsto
 \sum_{i,j,k}a_{ijk}x_i\otimes\mu_{w_j}(\tau)
                           \langle\delta(z_kG),v\rangle.
\]
Here $\delta(z_kG)$ are the three basis classes in
$H^1(C,\mathcal O_C)$ calculated in C.3ae, and the
pairing is Serre duality. This is a three-column map.
For the special triple $s_i=wz_i$, it sends the three
dual basis vectors to $x_i\otimes\mu_w(\tau)$ and
is injective. Therefore its rank-three locus is a
nonempty Zariski open. The map, its class $\tau$,
and its rank condition are defined over $K_0$ by
algebraic de Rham cohomology. It follows that our
algebraically independent triple satisfies this condition
as well as the transverse-intersection condition of
step 2. We use the special triple only to prove that
the rank-open set is nonempty; it is not asserted
to have the required finite base locus itself.

\emph{5. The obstruction on an actual source family.}
Let $P$ now be a smooth projective connected curve
and let $\Gamma\in\mathrm{CH}^2(X\times P)$.
The fixed closure on $\mathcal Z\times P$ used in
step 3 has refined restriction $b$ to $V\times P$.
The conormal sequence and the lifted Hodge Chern
character in C.3ae, with exactly the same rational
Chern-character and Gysin inputs, imply
\[
 \kappa\mathbin{\lrcorner}\operatorname{cl}_{\mathrm{Hdg}}(b)
       =0\quad\hbox{in }
 H^3(V\times P,L\otimes\Omega^1_{V\times P}),
 \qquad
 \kappa=\sum_{i,j,k}a_{ijk}x_iw_j\otimes\delta(z_kG).
\]
The first neighbourhood has this class because its
equation is $F+t\ell G$ and its only moving tangent
direction is that of $C$. No converse assertion
about deforming cycles is used.

Project to the Kunneth summand
\[
 H^0(\mathbb P^2,\mathcal O(1))\otimes
 H^1(W,\Omega_W^1(1))\otimes
 H^1(C,K_C)\otimes H^1(P,\mathcal O_P).
\]
The only contributing component of $b_0$ is
\[
 H^1(W,\Omega_W^1)\otimes H^0(C,K_C)
                                  \otimes H^1(P,\mathcal O_P).
\]
By the auxiliary Hodge calculation in step 1, its
$W$ factor lies in
$\operatorname{NS}(W)\otimes\mathbb C
 =\mathbb C H_W+\mathbb C\tau$.
This assertion applies uniformly to every $P$ and
every algebraic cycle $b_0$; it is not an assumption
that the other Hodge components are algebraic.
The $h b_1$ and $h^2b_2$ terms cannot contribute to
the displayed summand, since contraction acts only
on the $C$ differential and cannot remove their
projective-space differential. The $H_W$ component
is killed by each $\mu_{w_j}$. On the $\tau$
component the map is the injective three-column
map from step 4 tensored with $H^1(P,\mathcal O_P)$.
Hence that component is zero.

Pairing with $\tau$ now shows that
$\operatorname{pr}_{C\times P,*}(\tau b_0)$ has
zero mixed component in
$H^0(C,K_C)\otimes H^1(P,\mathcal O_P)$.
The two correction terms in step 3 do not affect
this conclusion. More explicitly, $q(W)=0$ gives
\[
       \operatorname{Pic}(W\times C\times P)
       =\operatorname{Pic}(W)\oplus
                         \operatorname{Pic}(C\times P).
\]
This follows from the Picard product description:
the possible cross homomorphisms involve
$\operatorname{Alb}(W)=\operatorname{Pic}^0(W)=0$.
If $b_1=D'+\lambda$ in this decomposition, projection
formula gives
\[
 \operatorname{pr}_{C\times P,*}(\tau a b_1)
                 =\deg(\tau D')\,\operatorname{pr}_C^*K_C.
\]
The other terms have either insufficient $W$
codimension or coefficient $\deg(\tau H_W)=0$.
Also $\operatorname{pr}_{C\times P,*}(\tau a^2b_2)=0$.
Thus the correction is constant in $P$, and the
line bundle $\rho_P(\Gamma|_{Y\times P})$ has zero
mixed Hodge component.

Subtracting its fibre at a base point gives a
morphism $P\to\operatorname{Pic}^0(C)$. By the
Jacobian universal property it factors through a
homomorphism $J(P)\to\operatorname{Pic}^0(C)$.
The vanishing mixed component is its zero differential.
Over $\mathbb C$ a homomorphism of abelian varieties
with zero differential is zero. The morphism is
therefore constant. Fibre compatibility of $\rho_P$
proves that $\rho$ kills every source-family difference.
Compactifying a smooth parameter curve and closing
its cycle in the fixed product reduces every ordinary
algebraic-equivalence generator to this case. Hence
\[
      \rho\bigl(\mathrm{CH}^2(X)_{\mathrm{alg}}|_Y\bigr)=0.
\]

\emph{6. A nonzero target and the integral cokernel.}
Pull back the divisors $D\subset W$ and a divisor
class on $C$ through $Y\to H\to W\times C$ and
take their product. This defines
\[
 \theta:\operatorname{Pic}^0(C)(\mathbb C)
                      \longrightarrow\mathrm{CH}^2(Y)_{\mathrm{alg}}.
\]
The Poincare family, restricted to smooth curves
joining the required points in its connected parameter,
proves ordinary algebraic triviality. The map is a
product of divisor pullbacks, so it uses no false
flatness assumption on $p$. Taking the same product
on $\mathcal D$ and using $p_*(h-2a)=1$ yields
\[
                              \rho\theta=-9\,\mathrm{id}.
\]
Multiplication by nine on the complex points of a
Jacobian is surjective. The asserted integral
cokernel surjection follows, and in particular the
ordinary rational restriction map is not surjective.

For a specified nontorsion class take
$o=(1:\eta:0)$ with $\eta^4=-1$ and
$p=(u:v:1)$ with $u$ transcendental over
$\overline{\mathbb Q}$ and $v^4=-1-u^4$.
The Abel closed immersion is defined over
$\overline{\mathbb Q}$ and every torsion point of
the Jacobian is algebraic over that field. Thus
$[p-o]$ is nontorsion, and $\theta([p-o])$ maps
to the nontorsion class $-9[p-o]$ in the cokernel.
All generic choices used for the surface and the
three sections have been specified separately; no
finite test of a Fermat surface is substituted for
the Hodge-generic surface in this proof.

The exact incidence and intersection arithmetic is retained in
\path{code/check_dimension_six_incidence.py}; its run took
$0.046509863$ seconds. The computation checks the displayed pushforward
polynomials, lattice pairings, coefficient counts and base degree.
The source-checked Hodge and genericity arguments above are separate
from those finite calculations.
